\chapter{Introducción}

En este capítulo se desarrollarán los siguientes temas:
\section*{Contenido}
\begin{enumerate}
\item ¿Qué es Simulación?
\item Clasificación de los Modelos de Simulación.
\item Un Primer Ejemplo.
\item Longitud de la Simulación.
\item Ventajas y Desventajas de los Modelos de Simulación.
\item Desarrollo de un Modelo de Simulación.
\item Resumen.
\end{enumerate}


\section{¿Qué es simulación?}

El Diccionario Básico Larousse define la palabra simulación como ``la acción de simular o fingir''. Posteriormente al buscar la palabra simular se dice que es ``dar la apariencia de lo que no se es, fingir'' o bien ``obtener la esencia de algo, sin ser real''.

En general, cualquier modelo que representa algo es un tipo de simulación. Es por lo tanto una palabra difícil de definir de forma apropiada. Para el desarrollo de la presente obra se entenderá el término de simulación como el desarrollo de un modelo lógico matemático de un sistema. Se busca obtener un modelo que logre una imitación del funcionamiento de un proceso de la vida real, o de un sistema virtual o físico, a través del tiempo.

De igual forma, un sistema se definirá como una colección de variables que interactúan entre sí dentro de ciertos límites para lograr un objetivo. Un modelo es una representación de los objetos del sistema y refleja de manera sencilla las actividades en las cuales se encuentran involucrados los objetos del sistema.

Desde este punto de vista la simulación consiste en construir modelos computacionales que describen las partes esenciales del comportamiento de un sistema, así como en diseñar y realizar experimentos con tales modelos con el fin de extraer conclusiones de sus resultados para apoyar la toma de decisiones. Se trata de una metodología de experimentación que se utiliza en el análisis de sistemas cuando las soluciones analíticas o matemáticas son muy complejas o no existen.




\section{Clasificación de los Modelos de Simulación.}


Los modelos con los cuales se puede experimentar en un sistema se suelen clasificar de acuerdo al siguiente esquema.
\begin{verbatim}

            Sistema
           /       \
          /         \
         /           \
Experimentar        Experimentar
con el              con un modelo
sistema             del sistema
                    /           \
                   /             \
                  /               \
              Modelo            Modelo
             físico           matemático
                             /          \
                            /            \
                           /              \
                       Solución        Simulación
                       analítica       por computadora

\end{verbatim}

En el presente curso estamos interesados en los modelos de simulación por computadora. Estos modelos se pueden subclasificar de diferentes formas, tal como se muestra a continuación.
\begin{itemize}
\item Estático o dinámico.

\item Determinístico o estocástico.

\item Discreto o continuo.

\item Centralizado o distribuido.
\end{itemize}

En el primer caso, el tiempo es el que decide de qué tipo de sistema se trata. Si éste cambia con el tiempo es dinámico. Si al transcurrir el tiempo el sistema no cambia es estático.

Los modelos determinísticos no tienen elementos basados en probabilidad, es decir no cuenta con procesos estocásticos en su estructura, por lo tanto su comportamiento se puede predecir. Los modelos estocásticos dependen de circunstancias que no se pueden predecir, dependen del resultado del azar y la aleatoriedad.

En un sistema de eventos discretos, el estado del sistema cambia sólo en ciertos puntos en el tiempo. Puede transcurrir una cierta cantidad de tiempo sin que cambie el estado del mismo. Por ejemplo, si se tiene una fila en un banco, puede transcurrir una cierta cantidad de tiempo sin que entre un nuevo cliente, o bien sin que ningún cliente sea atendido. En un sistema continuo, el estado del sistema cambia continuamente. Por ejemplo si se tiene una partícula en movimiento constante, cada instante de tiempo se encontrará en un lugar diferente.

El sistema puede ejecutarse en un solo lugar, o bien en distintos lugares que deben comunicarse entre sí para compartir la información producida.

Al final del presente capítulo, se presentan diferentes ejemplos para evaluar a cuál clasificación pertenecen. Sin embargo a lo largo de la presente obra estaremos interesados principalmente en sistemas dinámicos, estocásticos, discretos o continuos, que se pueden ejecutar en un único computador digital o de manera distribuida.



\section{Un Primer Ejemplo}

Se lanza una moneda legal tres veces. Supondremos que la moneda tiene de un lado una corona y del otro lado un escudo. Se desea obtener la probabilidad de obtener dos escudos.



\subsection{Solución con Teoría de Probabilidades}

De acuerdo con la teoría de probabilidades denominaremos a $x$ como una variable aleatoria que significa el número de escudos obtenidos. Se desea establecer el valor de $x=2$.

$$p(x=2) = \binom{3}{2}\times \frac{1}{2}^2 \times \frac{1}{2}^1$$  

$$p(x=2) = 3 \times 0.50^2 \times 0.50^1$$ 

$$p(x=2) = 0.375$$




\subsection{Solución con Simulación}


El problema anterior también podría resolverse mediante la técnica de simulación. Para ello se necesita utilizar un algoritmo que nos permita imitar el comportamiento del modelo. Todos los lenguajes de programación cuentan con una función que genera números aleatorios en el intervalo $[0,1]$. En realidad se trata de un generador de números pseudoaleatorios, pero hablaremos de eso más adelante.

Podemos usar el siguiente algoritmo para construir una moneda:

Procedimiento moneda
\begin{verbatim}
r := random(0,1)
Si 0<=r<=0.50 entonces es un escudo
en otro caso es una corona;
\end{verbatim}

Y podríamos repetir una gran cantidad de veces el proceso de tirar una moneda 3 veces, contando el número de veces que salen 2 escudos. Por ejemplo:
\begin{verbatim}
Repetir muchas veces
   Llamar a moneda 3 veces,
   Si ocurren 2 escudos contabilizar un éxito.
\end{verbatim}

La probabilidad de obtener dos escudos se puede calcular como:
\begin{verbatim}prob := número de éxitos / número de repeticiones.
\end{verbatim}

De esta forma, se obtendrá un valor cercano a 0.375


\section{La Distribución Binomial:}

Si se realiza el proceso anterior para calcular diferentes valores de "x", se obtendrá la distribución binomial dada por:
\begin{verbatim}
---------
x   p(x)
---------
0   0.125
1   0.375
2   0.375
3   0.125
---------
\end{verbatim}



\section{Longitud de la Simulación}


Un sistema en el que existen puntos de inicio y puntos de terminación precisos se conoce como un sistema con punto de finalización. Por ejemplo si se simula un sistema de lotería es claro que el sistema finaliza cuando se han sacado todos los números ganadores.

Existen otros sistemas que siempre están en curso y que carecen de un punto de inicio y terminación precisos. Estos sistemas se conocen como sistemas sin punto de finalización. Por ejemplo el proceso de producción de una fábrica siempre está en marcha, nunca se detiene, las condiciones al final de un día son las condiciones iniciales del día siguiente.

Cuando un sistema tiene un punto de finalización, la longitud de la simulación, la cantidad de tiempo sobre la cual se debe conducir el experimento para obtener los análisis deseados va desde su punto de inicio hasta su punto de terminación. Para obtener resultados confiables se debe determinar cuantas veces se debe repetir la simulación.

En contraste para un sistema de no terminación, se debe elegir no solo el número de veces que se va a repetir la simulación, sino también los puntos iniciales y finales de la misma. Por ejemplo, para una fábrica se podrían establecer ciertas condiciones iniciales y se simula el proceso de producción durante un día, y luego se repite ese proceso 30 veces para simular un mes.

Al igual que los problemas de condiciones iniciales y de terminación, existen otros problemas tácticos asociados al desarrollo de un modelo de simulación. Muchos estudios de simulación se refieren a sistemas de investigación que operan en una condición de estado estable. Por desgracia un sistema de simulación no inicia necesariamente en un estado estable. Debe arrancarse y detenerse. Debido a ciertas condiciones artificiales que se introducen, al iniciar de manera brusca las operaciones, el desempeño de los sistemas simulados no se vuelve representativo del sistema real correspondiente hasta que alcanzado una condición de estado estacionario o estable. El estado estable se alcanza cuando las distribuciones de probabilidad del sistema han alcanzado su punto de equilibrio unas entre otras. El tiempo que tarda un sistema en alcanzar este punto varía de sistema a sistema y de simulación a simulación.



\section{Ventajas y Desventajas de los Modelos de Simulación.}

Hemos definido la simulación como la experimentación con un modelo de un sistema real o virtual. Si el sistema no existe, es evidente la necesidad de estudiar su comportamiento antes de construirlo. Ahora bien si el sistema existe en la realidad, el estudio del mismo puede ocasionar problemas en el sistema en funcionamiento. Por ejemplo:
\begin{itemize}

\item Puede ocasionar problemas en las operaciones cotidianas.

\item Si hay personas involucradas en el sistema puede ocurrir el llamado \textit{efecto Hawthorne}, el cual establece que cuando la gente sabe que es observada puede oficiar su comportamiento.

\item Es difícil reproducir las mismas condiciones de funcionamiento para repetir un experimento.

\item Puede consumir más tiempo y ser más costoso que construir un modelo y probar con el modelo.

\item No es posible experimentar con alternativas radicalmente diferentes a la forma como el modelo funciona.
\end{itemize}

Por esta razón, se debe considerar ventajoso utilizar un modelo de simulación cuando se dan algunas de las siguientes condiciones:

\begin{itemize}
\item No existe una formulación matemática completa del modelo o bien no se han desarrollado métodos analíticos para trabajar con el modelo.

\item Los métodos matemáticos existen, pero son tan complejos y arduos que la simulación puede ser un modelo más sencillo.

\item Los sistemas matemáticos o analíticos deben hacer gran número de suposiciones para mantenerlos sencillos, mientras que en los modelos de simulación se pueden enfrentar problemas más complejos.

\item Existen modelos matemáticos, pero no se cuenta con el personal calificado para desarrollar estos modelos. En este caso un modelo de simulación es mucho más factible de ser desarrollado.

\item La simulación puede ser la única posibilidad, por la imposibilidad o el peligro de conducir experimentos en el modelo real.

\item Algunos sistemas requieren de largos períodos de tiempo para producir resultados. En los sistemas de simulación digital se puede "comprimir el tiempo" para obtener resultados con mayor rapidez.

\item Una vez construido el modelo puede ser modificado de manera relativamente rápida para analizar diferentes esquemas.

\subitem Los sistemas de simulación incluyen gráficos e información que resulta más fácil de comprender que los métodos matemáticos o analíticos.

\end{itemize}

Algunas ventajas adicionales de la simulación son sus efectos de educación y entrenamiento. El desarrollo y el uso de uno de estos sistemas le permiten al experimentador sentir y jugar con el sistema de manera tal que se puede adquirir un conocimiento más profundo de su comportamiento.


Entre sus desventajas se pueden citar:
\begin{itemize}
\item Los modelos de simulación por computadora son muy costosos y requieren tiempo para su desarrollo y validación.

\item Se requiere gran cantidad de ejecuciones computacionales para obtener soluciones óptimas. Muchas veces solo es posible encontrar soluciones aproximadas.

\item Los resultados de los modelos pueden ser imprecisos, lo que dificulta la toma de decisiones.

\item Los modelos de simulación pueden dar un falso sentido de seguridad, para evitar esto deben de planearse y validarse fuertemente.

\item Muchas organizaciones no están acostumbradas a trabajar con estos modelos por lo que no tienen confianza en sus resultados.
\end{itemize}



\section{ Desarrollo de un Modelo de Simulación}

La gran mayoría de las metodología para el desarrollo de sistemas de simulación incluye los siguientes pasos:


\subsection{Definición del Sistema}
Todo estudio inicia con una definición del sistema que se desea simular. Se deben tomar en cuanta los siguientes elementos: 
\begin{itemize}
\item Objetivo del sistema.
\item Variables a considerar.
\item Medidas de optimalidad.
\item Variables no controladas
\item Intervalos de confianza estadísticos.
\end{itemize}

Deben describirse las interacciones entre las variables de decisión para lograr optimizar el objetivo. Surgen además algunas definiciones.

Se conoce como "estado del sistema" el valor de las variables en un instante dado.

Un ``evento'' representa un acontecimiento que cambia el valor del estado del sistema.


\subsection{Recolección de Datos}


Es posible que la facilidad o dificultad para obtener los datos influya sobre la formulación del modelo. Además de obtener los datos se deben realizar pruebas estadísticas para valorar la confiabilidad de los mismos.


\subsection{Formulación del Modelo}

Una vez que se han definido los resultados que se esperan obtener se debe construir el modelo. Este corresponde al modelo lógico. 


\subsection{Implementación del Modelo}

Con el modelo lógico definido se debe seleccionar un lenguaje computacional para implementarlo. Se pueden utilizar lenguajes de programación de uso general o bien lenguajes de programación especiales para simulación.


\subsection{Validación}


Esta es una de las etapas más importantes. A través de esta etapa se prueba que lo que se ha construido responde al modelo original. Algunas formas para validar un modelo son:
\begin{itemize}
\item Opinión de expertos sobre el comportamiento del mismo

\item Utilizar datos históricos que sirvan para valorar la predicción del modelo.

\item Utilizar el modelo para producir datos futuros y evaluar su resultado.
\end{itemize}

\subsection{Experimentación}


La experimentación se realiza una vez que el modelo se ha validado. En esta etapa se trata de establecer la sensibilidad del modelo ante variantes en las variables de entrada y variables del sistema.


\subsection{Documentación}

Como en todo sistema computacional se deben desarrollar las adecuadas documentaciones del programa, manuales de usuario, ayuda en línea, etc.


Una vez que se ha terminado con las etapas de construcción del modelo se puede utilizar para la toma de decisiones y para experimentar con cambios en el sistema real.



\section{Resumen}


La simulación busca construir un modelo artificial que se comporte como el modelo al cual se quiere imitar. 

Antes de proceder a construir un simulación de un modelo se deben tomar en consideración las ventajas y desventajas del uso de este enfoque.

Siempre y cuando el problema se pueda resolver mediante un modelo analítico se debe utilizar esta solución. En caso que no existe un modelo analítico, o este sea muy complejo, se debe utilizar la simulación.

Los etapas para construir un modelo de simulación son:
\begin{itemize}
\item Definición del Sistema
\item Recolección de Datos
\item Formulación del Modelo
\item Implementación del Modelo.
\item Validación.
\item Experimentación.
\item Documentación
\end{itemize}


 
\section{Ejercicios}

\begin{enumerate}
    \renewcommand{\theenumi}{Ejercicio \arabic{enumi}} % Cambia el formato de numeración
    \setcounter{enumi}{0} % Inicia la numeración desde 1
    \item Vaya al sitio www.statistics101.net e instale el lenguaje de simulación que se utilizará en el curso.
    \item Se tienen 3 monedas. Las cuales pueden dar como resultado escudo o corona.
        \begin{enumerate}[label=\alph*.]
        \item Utilizando probabilidad matemática calcule el valor de obtener dos coronas.
        \item Con un programa de simulación, calcule la probabilidad de obtener 0 coronas.
        \item Con un programa de simulación, calcule la probabilidad de obtener 1 coronas.
        \item Con un programa de simulación, calcule la probabilidad de obtener 2 coronas.
        \item Con un programa de simulación, calcule la probabilidad de obtener 3 coronas
        \item Qué clase de distribución de probabilidad es esta?
      \end{enumerate}

Sugerencia:
\subitem Probabilidad de obtener 0 coronas: 0.125
\subitem Probabilidad de obtener 1 corona:  0.375
\subitem Probabilidad de obtener 2 coronas: 0.375
\subitem Probabilidad de obtener 3 coronas: 0.125

    \item Construya una distribución binomial con un máximo de 3 éxitos y una probabilidad de obtener un éxito de 0.50. Grafique un histograma con los resultados obtenidos   
    \item Se tiene un mazo de cartas convencional.
    \begin{enumerate}[label=\alph*.]
        \item ¿Cuál es la probabilidad de sacar dos cartas y que ambos sean ases?
        \item Construya un programa de simulación que genere estos resultados.

Sugerencia: La probabilidad es: $\frac{1}{221}$ $(0.0045)$
      \end{enumerate}
    \item Se tiene un mazo de cartas convencional.
        \begin{enumerate}[label=\alph*.]
        \item ¿Cuál es la probabilidad que al sacar una carta se obtenga una espada o un trébol?
        \item Construya un programa de simulación que genere estos resultados.

      \end{enumerate}
Sugerencia: La probabilidad es: 0.50


    \item Se tiene un mazo de cartas convencional, el cual ha sido barajdo. Se reparten 5 cartas y se desea establecer la probabilidad que dentro de esas 5 cartas se encuentre un par de ases.
        \begin{enumerate}[label=\alph*.]
        \item Resuelva el problema utilizando probabilidad clásica. 
        \item Resuelva el problema construyendo un programa de simulación
      \end{enumerate}

    \item Se tiene un mazo de cartas convencional, el cual ha sido barajdo. Se reparten 5 cartas y se desea establecer la probabilidad que las 5 cartas sumen entre 17 y 21.\\
Resuelva el problema construyendo un programa de simulación.
    \item Se tiene un juego con una probabilidad de ganar de 0.48 y una probabilidad de perder de 0.52. Cada vez que se apuesta una cantidad de dinero se gana lo que se apuesta o se pierde todo. Se inicia el juego con \$100 dólares y se desea llegar a obtener \$200.

\subitem Si se utiliza una estrategia de apuesta $(x,x)$ con $x=10$. Es decir, si gano apuesto $x$ y si pierdo apuesto $x$, determine la probabilidad de obtener la cantidad deseada.

\subitem Si se utiliza una estrategia de apuesta $(x,2\times x)$ con $x=10$. Es decir, si gano apuesto $x$, pero si pierdo apuesto el doble de lo que perdí $2\times x$, determine la probabilidad de obtener la cantidad deseada.

\end{enumerate}



\subsection{Ejercicio avanzado}

La Agencia Espacial Internacional desea mejorar su servicio de comunicación por satélite de Tokio a Florida. En el sistema actual, los mensajes pueden enviarse de Tokio hacia Florida utilizando la red de satélites que se muestra a continuación.

\begin{verbatim}
       .----> S1 ---.
      /              \
     /0.90            \0.80
    /                  \
   /                    '--->
Tokio                         Florida
   \                    .--->
    \                  /
     \0.90            /0.70
      \              /
       '----> S2 ---'
\end{verbatim}


La ruta superior implica enviar un mensaje de Tokio a Florida por medio del satélite S1. La segunda alternativa es enviar un mensaje de Tokio a Florida por medio del satélite S2. Debido a las órbitas actuales estos satélites pueden conectarse entre sí y entre las ciudades solo una fracción del tiempo. Estas fracciones se escriben junto a cada nodo de la red. Por ejemplo el valor de 0.90 que une a Tokio con S1 significa que esta conexión tiene un 90\% de probabilidad de que un mensaje enviado de Tokio llegue a S1.

\begin{enumerate}
\item Determine utilizando probabilidades, cuál es la probabilidad total que un mensaje enviado desde Tokio llegue a Florida, dada la red actual de satélites. Es decir, la probabilidad para el evento $T \rightarrow F$, dada por$ ( P(T \rightarrow F) )$.


\item Se desea mejorar al nivel actual de comunicaciones, para ello se están evaluando algunas posibilidades. El personal ha descubierto dos opciones para tratar de mejorar el actual sistema de satélites. La primera consiste en alterar la órbita de S2, lo que aumentaría las probabilidades de comunicarse con Florida a un nuevo valor de  0.80.


\item Con un esfuerzo sustancialmente mayor, se puede modificar la órbita de S1 para que S1 y S2 puedan comunicarse hacia adelante y hacia atrás con una probabilidad de 0.80. Bajo estas condiciones, ¿qué tan fácil es calcular las nuevas probabilidades?

\end{enumerate}
\begin{verbatim}
       .----> S1 ---.
      /       |      \
     /0.90    |       \0.80
    /         |        \
   /          |         '--->
Tokio         |0.80           Florida
   \          |         .--->
    \         |        /
     \0.90    |       /0.70
      \       |      /
       '----> S2 ---'
\end{verbatim}